\chapter{Auditing Compatibility with xna4-spec}
\index{xna4-spec}\index{API audit}
\label{ch:xna4-spec-auditing}

\texttt{xna4-spec} is an offline, machine-readable conversion of Microsoft's XNA~4.0
documentation.  It is useful for asking what the documented API surface contains.  It is not part
of CNA's build, test, or oracle chain, and it is not the source of CNA's dated 227-of-245 FNA type
count; that number was produced by diffing FNA's public type list against CNA headers.

\section{What the repository contains}

At pinned commit \texttt{8f61207}, the top-level index references 544 XML type files across 19
documented namespaces.  Each file records a kind, syntax, members, overloads, parameters, return
text, remarks, and links derived from the former Microsoft documentation archive.  The collection
includes runtime namespaces and the content-pipeline namespaces CNA does not aim to implement.

The README's example tree says ``550 types,'' while its coverage table, index, plan numbering, and
files all say 544.  This small internal drift is the first rule of using the corpus: treat even a
machine-readable conversion as data to validate, not as an infallible executable specification.
The repository's disclaimer also makes clear that the converted documentation remains Microsoft's
content and is offered for informational and educational use.

\section{Where it sits in CNA's authority hierarchy}

No file in CNA references \texttt{xna4-spec}.  CNA's current port checklist instead names the XNA
reference-assembly XML as the authority for API surface and FNA as the reference for behavior.
Real XNA execution, FNA differential values, and image oracles then settle questions documentation
cannot.

\texttt{xna4-spec} is therefore a convenient complementary inventory.  It can expose a missing
property that FNA itself omitted, supply remarks for manual review, and make namespace-wide audits
scriptable.  It cannot prove CNA behavior, and a clean comparison against it is not an existing CNA
gate.  A book or report must not promote an offline sibling repository into a runtime authority
merely because its name contains ``spec.''

\section{A reproducible surface audit}

For one type, a disciplined comparison has six steps:

\begin{enumerate}
  \item pin the \texttt{xna4-spec}, CNA, FNA, and reference-assembly revisions;
  \item extract fields, properties, constructors, methods, events, and overload signatures from
        the XML rather than counting prose mentions;
  \item map C\# idioms explicitly: properties to getter/setter names, operators to C\texttt{++}
        operators, \texttt{ref}/\texttt{out} to reference overloads, and iterators to CNAEXT;
  \item compare CNA's public header and classify each row as exact surface, accepted language
        adaptation, marked extension, unexplained extra, or missing;
  \item read remarks and implementation for every hand-written accessor or behavior-bearing
        method;
  \item attach an executable oracle for semantics and record cases the XML cannot decide.
\end{enumerate}

The mapping table is part of the result.  Otherwise \texttt{Equals(object)} looks like a missing
member even though a C\texttt{++} value type has no universal boxed root, and a default constructor
or iterator can look like an unmarked XNA extension even when it exists only for the target
language's ordinary value/container idiom.

\section{Worked example: Ray}

\texttt{Microsoft.Xna.Framework/Ray.xml} describes the \texttt{Position} and \texttt{Direction}
fields, the two-vector constructor, equality, hash and string methods, intersection overloads, and
equality operators.  CNA's \texttt{modules/math/.../Ray.hpp} contains both fields, the constructor,
value equality, \texttt{ToString}, four intersection target types, output-reference forms for box,
sphere and plane, and free equality operators.

The comparison also exposes adaptations.  CNA adds a default constructor, omits boxed
\texttt{Equals(Object)}, and returns \texttt{std::size\_t} from \texttt{GetHashCode} rather than
C\# \texttt{int}.  Those choices match CNA's documented C\texttt{++} conventions; they should be
classified, not silently counted as either perfect matches or defects.

The converted XML itself needs normalization.  Its \texttt{Intersects} list contains duplicate
output overload records, and some signatures are labelled with the wrong target type even though
the surrounding parameter text names Plane or BoundingSphere.  A script that compares signature
strings blindly will invent duplicates and gaps.  Cross-checking the reference-assembly member IDs
or the original documentation is necessary whenever converted fields disagree.

\section{Worked example: Viewport}

The Viewport XML lists two non-default constructors, nine properties, and
\texttt{Project}, \texttt{Unproject}, and \texttt{ToString}.  CNA's header provides those members
through a mix of generated field properties and hand-written accessors.  It also has a default
constructor that produces zero width/height and a 0-to-1 depth range.

The surface comparison looks strong, but \texttt{TitleSafeArea} demonstrates why surface is not
semantics.  CNA returns \texttt{Bounds} unchanged and pins that choice with
\texttt{TitleSafeAreaEqualsBounds}.  The documentation describes a subset guaranteed visible on
lower-quality displays.  A passing test proves CNA's chosen behavior is stable; it does not prove
the choice matches XNA on every platform.

The right outcome is a recorded behavioral deviation or a real-XNA experiment, not a surface-
coverage percentage.  Hand-written properties deserve particular scrutiny because computation and
policy can hide behind a signature that compares perfectly.

\section{What inventories cannot prove}

An XML API inventory cannot establish default values produced by a runtime, floating-point edge
behavior, exception type and timing, lifetime, GPU state, wire compatibility, content decoding, or
platform substitution.  It also cannot tell whether a nominally present method is a stub, a safe
refusal, or a complete implementation.

Use \texttt{xna4-spec} to generate questions and completeness rows.  Use CNA's compile-time CNAEXT
gate to detect unmarked additions, reference assemblies to settle surface, FNA or real XNA to
settle semantics, and the oracle hierarchy from Chapter~\ref{ch:oracles} to produce verdicts.  The
most trustworthy audit is not the one with the highest percentage, but the one that preserves
these authorities instead of collapsing them into one number.
